\documentclass[conference]{IEEEtran}

\usepackage[utf8]{inputenc}
\usepackage[T1]{fontenc}
\usepackage[brazil]{babel}
\usepackage{amsmath,amssymb,amsfonts}
\usepackage{graphicx}
\usepackage{booktabs}
\usepackage{array}
\usepackage{cite}
\usepackage{float}
\usepackage{algorithm}
\usepackage{algpseudocode}

\begin{document}
\title{Otimização do Custo e Qualidade de Pilhas de Minério - Entrega \#1}

\author{
\IEEEauthorblockN{Saulo Bruno do Amaral - 2024425199}
\IEEEauthorblockN{Letícia Braga Esteves - 2023422609}
\IEEEauthorblockN{Yuri - xxx}
\IEEEauthorblockN{Roberto - xxx}
\IEEEauthorblockN{Juliana - xxx}
\IEEEauthorblockN{Igor - xxx}\IEEEauthorblockA{
Universidade Federal de Minas Gerais -- UFMG\\
Belo Horizonte, Brasil
}
}

\maketitle

\section{Introdução}

O problema consiste na composição de pilhas de minério destinadas a dois equipamentos de sinterização. Cada pilha deve atingir uma massa final definida, respeitar limites de qualidade e utilizar apenas minérios compatíveis com o equipamento ao qual está associada. A decisão é realizada em unidades discretas de caminhões, cada um com capacidade fixa de 2 kt.

Nesta primeira entrega, o problema é tratado de forma mono-objetivo. São consideradas separadamente três funções objetivo: minimização do custo total de aquisição dos minérios, minimização do desvio quadrático do teor de SiO$_2$ em relação ao valor-alvo e minimização do desvio quadrático do teor de Al$_2$O$_3$ em relação ao valor-alvo.

A solução computacional é construída a partir de uma representação discreta das pilhas...

\section{Formulação Matemática}

\subsection{Conjuntos e Índices}

\begin{align}
I &= \{1,\ldots,|I|\}, \\
J &= \{1,\ldots,|J|\}, \\
S &= \{1,\ldots,|S|\}, \\
K &= \{1,2\}.
\end{align}

\begin{align}
k=1 &\rightarrow \mathrm{SiO_2}, \\
k=2 &\rightarrow \mathrm{Al_2O_3}.
\end{align}

\subsection{Parâmetros}

\begin{table}[htbp]
\caption{Parâmetros do problema}
\centering
\begin{tabular}{p{0.20\columnwidth} p{0.70\columnwidth}}
\toprule
\textbf{Símbolo} & \textbf{Parâmetro} \\
\midrule
$P_j$ & Preço do minério $j$ \\
$F_{kj}$ & Teor da substância $k$ no minério $j$ \\
$D_j^{\min}$ & Disponibilidade mínima do minério $j$ \\
$D_j^{\max}$ & Disponibilidade máxima do minério $j$ \\
$M_i$ & Massa da pilha $i$ \\
$T$ & Capacidade de um caminhão \\
$B_{sj}$ & Compatibilidade entre o sinter $s$ e o minério $j$ \\
$L_{ks}^{\min}$ & Limite mínimo da substância $k$ no sinter $s$ \\
$L_{ks}^{\max}$ & Limite máximo da substância $k$ no sinter $s$ \\
$L_{ks}^{\mathrm{alvo}}$ & Valor-alvo da substância $k$ no sinter $s$ \\
$s(i)$ & Sinter associado à pilha $i$ \\
\bottomrule
\end{tabular}
\label{tab:parametros}
\end{table}

\subsection{Variáveis de Decisão}

Seja $x_{ij}$ a quantidade de caminhões contendo o minério $j$ destinados à pilha $i$.

\begin{equation}
 x_{ij} \in \mathbb{Z}_{\geq 0},
 \qquad \forall i\in I,\\ \forall j\in J
\end{equation}

\subsection{Teor das Pilhas}

\begin{equation}
L_{ki}
=
\frac{T}{M_i}
\sum_{j\in J} x_{ij} F_{kj}
\end{equation}

\subsection{Funções Objetivo}

\subsubsection{Minimização do custo}

\begin{equation}
\min f_1(x)
=
T
\sum_{i\in I}
\sum_{j\in J}
P_j x_{ij}
\end{equation}

\subsubsection{Minimização do desvio de SiO$_2$}

\begin{equation}
\min f_2(x)
=
\sum_{i\in I}
\left(
L_{1i}
-
L_{1,s(i)}^{\mathrm{alvo}}
\right)^2
\end{equation}

\subsubsection{Minimização do desvio de Al$_2$O$_3$}

\begin{equation}
\min f_3(x)
=
\sum_{i\in I}
\left(
L_{2i}
-
L_{2,s(i)}^{\mathrm{alvo}}
\right)^2
\end{equation}

\subsection{Restrições}

\subsubsection{Massa das pilhas}

\begin{equation}
T
\sum_{j\in J}
x_{ij}
=
M_i,
\qquad \forall i\in I
\end{equation}

\subsubsection{Qualidade}

\begin{equation}
L_{k,s(i)}^{\min}
\leq
L_{ki}
\leq
L_{k,s(i)}^{\max},
\qquad
\forall i\in I,\ \forall k\in K
\end{equation}

\subsubsection{Compatibilidade}

\begin{equation}
x_{ij}
\leq
\frac{M_i}{T}
B_{s(i),j},
\qquad
\forall i\in I,\ \forall j\in J
\end{equation}

\subsubsection{Disponibilidade}

\begin{equation}
D_j^{\min}
\leq
\sum_{i\in I} x_{ij}
\leq
D_j^{\max},
\qquad
\forall j\in J
\end{equation}

\subsubsection{Domínio}

\begin{equation}
x_{ij} \in \mathbb{Z}_{\geq 0},
\qquad
\forall i\in I,\ \forall j\in J
\end{equation}

\section{Representação Computacional da Solução}

\subsection{Representação de uma Pilha}

\subsection{Representação da Solução Completa}

\subsection{Relação com $x_{ij}$}

\subsection{Verificação de Viabilidade}

\section{Heurística Construtiva}

\subsection{Geração da Solução Inicial}

\subsection{Atendimento das Disponibilidades Mínimas}

\subsection{Preenchimento das Pilhas}

\subsection{Tratamento de Inviabilidade Inicial}

\section{Metaheurística VNS/GVNS}

\subsection{Estrutura Geral}

\subsection{Vizinhança $N_1$}

\subsection{Vizinhança $N_2$}

\subsection{Vizinhança $N_3$}

\subsection{Busca Local}

\subsection{Perturbação}

\subsection{Tratamento de Soluções Inviáveis}

\subsection{Critério de Aceitação}

\subsection{Critério de Parada}

\subsection{Pseudocódigo}

\begin{algorithm}[H]
\caption{Metaheurística VNS/GVNS}
\label{alg:gvns}
\begin{algorithmic}[1]
    \State ...
\end{algorithmic}
\end{algorithm}

\section{Planejamento Experimental}

\subsection{Configuração dos Experimentos}

\subsection{Parâmetros do Algoritmo}

\subsection{Número de Execuções}

\subsection{Métricas Utilizadas}

\subsection{Critério de Comparação}

\section{Resultados e Discussões}

\subsection{Resultados para $f_1$}

\subsection{Resultados para $f_2$}

\subsection{Resultados para $f_3$}

\subsection{Curvas de Convergência}

\subsection{Melhores Soluções Encontradas}

\section{Conclusão}

\end{document}
